\documentclass[11pt,a4paper]{article}
\usepackage[margin=25mm]{geometry}
\usepackage[T1]{fontenc}
\usepackage{lmodern}
\usepackage{amsmath,amssymb,amsthm,mathtools}
\usepackage[hidelinks]{hyperref}
\usepackage{microtype}
\newtheorem{theorem}{Theorem}
\newtheorem{lemma}[theorem]{Lemma}
\newtheorem{corollary}[theorem]{Corollary}
\newtheorem{remark}[theorem]{Remark}
\newtheorem{proposition}[theorem]{Proposition}
\newcommand{\N}{\mathbb N}
\newcommand{\one}{\mathbf 1}
\title{A Sharp Obstruction to Proving Odd-Core Tightness from Dyadic Flux Alone}
\author{Independent research note --- for mathematical review}
\date{17 November 2026}
\begin{document}
\maketitle
\begin{abstract}
The odd-core compression result in our ongoing study reduces a clipping-tail
problem for $3x+c$ to uniform integrability of the stopped logarithmic
predecessor potential on \emph{odd} targets.  The recurrence
$h(y)-h(2y)=h((2y-c)/3)$ suggests a dyadic telescoping argument.
We show that such an argument cannot succeed using harmonicity,
positivity, cycle-level conservation, a bounded weighted Ces\`aro first
moment, and eventual entrance into a fixed cycle alone.  For the
$3x+1$ map we construct an explicit $\{0,1\}$-valued harmonic profile
supported on one lacunary inverse orbit, with uniformly bounded weighted
Ces\`aro mass but arbitrarily high odd-target spikes.  All statements
are about an \emph{abstract harmonic profile}; this is not the actual
predecessor density.  The construction identifies the genuinely missing
input: quantitative information about how logarithmic predecessor mass
splits among inverse branches.
\end{abstract}

\section{Notation and the proposed dyadic route}
For $n\in\N$ let
\[
 T(n)=\begin{cases}n/2&n\text{ even},\\(3n+1)/2&n\text{ odd}.
 \end{cases}
\]
For $y\in\N$ let $q(y)=(2y-1)/3$ when this is a positive odd integer,
and otherwise regard $q(y)$ as absent.  A formal harmonic predecessor
profile is a function $H:\N\to[0,1]$ satisfying
\begin{equation}\label{eq:recurrence}
 H(y)=H(2y)+\one_{\{q(y)\text{ exists}\}}H(q(y))
 \qquad(y\notin\{1,2\}).
\end{equation}
The cycle $\{1,2\}$ is excluded, as required in the predecessor-potential
recurrence used in the consolidated manuscript.  Actual logarithmic
predecessor densities satisfy this recurrence; the converse is not
asserted.  Given $0<\beta<1$, write $W_R=(R^{1-\beta},R]\cap\N$.
We ask whether \eqref{eq:recurrence}, together with bounded weighted
Ces\`aro first moments, forces
\begin{equation}\label{eq:odd-ui}
 \lim_{M\to\infty}\limsup_{R\to\infty}\frac1R
 \sum_{\substack{a\in W_R\\a\ \mathrm{odd}\\aH(a)>M}}aH(a)=0.
\end{equation}
The answer is \emph{no}, even under additional natural restrictions.

\section{An explicit odd harmonic spine}
The choice of large, but \emph{uniformly bounded}, inverse Syracuse
valuations is essential to keeping the example elementary.
\begin{lemma}[An arithmetic inverse branch]\label{lem:branch}
For every positive odd $u$ with $3\nmid u$ there is a unique
$e(u)\in\{8,9,10,11,12,13\}$ satisfying
\begin{equation}\label{eq:cong}
 2^{e(u)}u\equiv4\pmod 9.
\end{equation}
The number
\begin{equation}\label{eq:inverse}
 v=\frac{2^{e(u)}u-1}{3}
\end{equation}
is a positive odd integer, $v\equiv1\pmod3$, $v\ge85u$, and
$T(v)=2^{e(u)-1}u$.
\end{lemma}
\begin{proof}
The element $2$ has multiplicative order $6$ modulo $9$ and hence
$2^8,\ldots,2^{13}$ run through all six units modulo $9$.
Because $u$ is a unit, precisely one of these powers satisfies
\eqref{eq:cong}.  The congruence implies $2^{e(u)}u-1\equiv3\pmod9$,
so $v$ is integral, is $1$ modulo $3$, and is odd since the numerator
is odd.  As $e(u)\ge8$ and $u\ge1$,
$v\ge(256u-1)/3\ge85u$.  The last identity follows from
$3v+1=2^{e(u)}u$ and the odd branch of $T$.
\end{proof}

Define $u_0=5$ and, recursively,
\[
 e_i=e(u_i),\qquad u_{i+1}=\frac{2^{e_i}u_i-1}{3}.
\]
By Lemma~\ref{lem:branch}, every $u_i$ is odd, coprime to $3$, and
$u_{i+1}\ge85u_i$.  Define the subset
\begin{equation}\label{eq:support}
 \mathcal S=\{1,2,4,8\}\ \cup\
 \bigcup_{i\ge0}\{2^k u_i:0\le k<e_i\},
 \qquad H(n)=\one_{\{n\in\mathcal S\}}.
\end{equation}
For example, $u_0=5$, $e_0=9$, $u_1=853$, $e_1=10$,
$u_2=291157$.  Each inverse segment is
\[
 u_{i+1}\ \xrightarrow{\ T\ }\ 2^{e_i-1}u_i
 \ \xrightarrow{\ T\ }\ 2^{e_i-2}u_i
 \ \xrightarrow{\ T\ }\ \cdots\ \xrightarrow{\ T\ }\ u_i.
\]
The initial segment is $5\to8\to4\to2\to1\to2$.

\begin{theorem}[Bounded mean, unbounded odd spikes]\label{thm:spine}
The profile $H$ of \eqref{eq:support} has all the following properties.
\begin{enumerate}
\item $0\le H\le1$, $H(1)=H(2)=1$, and $H$ satisfies the full
      nonperiodic recurrence \eqref{eq:recurrence}.
\item Every state in the support of $H$ reaches the $1$--$2$ cycle.
      All states in the support are coprime to $3$.
\item On every backward-tree level outside the $1$--$2$ cycle,
      the sum of $H$ is exactly $1$.
\item For every real $R\ge1$,
      \begin{equation}\label{eq:mean}
       \frac1R\sum_{n\le R}nH(n)\le10.
      \end{equation}
\item For every $0<\beta<1$, condition \eqref{eq:odd-ui} fails;
      more precisely,
      \begin{equation}\label{eq:odd-spike}
       \lim_{M\to\infty}\limsup_{R\to\infty}
       \frac1R\sum_{\substack{a\in W_R\\ a\ \mathrm{odd}\\aH(a)>M}}
       aH(a)\ge1.
      \end{equation}
      For every $p>1$, the odd-target $L^p$ means
      $R^{-1}\sum_{a\le R,\,a\ \mathrm{odd}}(aH(a))^p$ are unbounded.
\end{enumerate}
\end{theorem}
\begin{proof}
The displayed inverse segments concatenate to a single injective
infinite backward ray emanating from the cycle.  Its supported
vertices all map forward to supported vertices.  Outside the cycle,
every supported vertex has \emph{exactly one} supported immediate
predecessor: the next doubled vertex in its own block, or the next
odd vertex $u_{i+1}$ at a block's upper endpoint; the same holds for
$4$ and $8$.  No unsupported vertex has a supported predecessor,
since the support is forward invariant.  Consequently the indicator
$H$ satisfies \eqref{eq:recurrence}, with the periodic targets $1,2$
properly excluded.  Lemma~\ref{lem:branch} gives coprimality with $3$;
forward convergence and cycle-level conservation follow from the
single-ray construction.

For the weighted first moment, if $R<5$ only the finite initial segment
contributes and \eqref{eq:mean} is immediate.  Otherwise choose the
unique $i\ge0$ with $u_i\le R<u_{i+1}$.  The initial segment has
weight $15\le3R$.  The contribution from the current block $i$,
truncated at $R$, is a geometric series of powers of two and is at
most $2R$.  Every earlier block $j<i$ has total weight
\[
 \sum_{k=0}^{e_j-1}2^ku_j
 =(2^{e_j}-1)u_j=3u_{j+1}+1-u_j\le3u_{j+1}+1.
\]
Since $u_{j+1}\ge85u_j$, one has
$\sum_{\ell=1}^iu_\ell\le(85/84)u_i$ and $i\le u_i$.  Thus
\[
 \sum_{n\le R}nH(n)
 \le 3R+2R+3(85/84)u_i+i
 \le(6+255/84)R<10R.
\]
Here over-counting entire earlier blocks is harmless even if their
largest vertices exceed $R$.

Finally, take $R=u_i\to\infty$.  The number $u_i$ is odd, belongs
to $W_{u_i}$, and satisfies $H(u_i)=1$.  For every fixed $M$ this
single term contributes $u_i/R=1$ to \eqref{eq:odd-spike} for large
$i$.  Similarly, its contribution to the normalized $p$th moment
is $u_i^{p-1}\to\infty$.
\end{proof}

\section{Why dyadic telescoping does not smooth weighted spikes}
For an odd base $u$ write $H_k=H(2^ku)$.  When $2^ku$ is not
periodic, the recurrence yields
\[
 H_k-H_{k+1}=
 \one_{\{(2^{k+1}u-1)/3\ \text{positive odd integer}\}}
 H\left(\frac{2^{k+1}u-1}{3}\right)\ge0.
\]
The unweighted differences telescope, but the odd-core moment uses a
factor comparable to $2^ku$.  For any nonnegative weights $w_k$,
Abel summation gives the exact identity
\begin{equation}\label{eq:abel}
 \sum_{k=K}^M w_k(H_k-H_{k+1})
 =w_K H_K+\sum_{k=K+1}^M(w_k-w_{k-1})H_k-w_MH_{M+1}.
\end{equation}
At the branching edge of our construction,
$H(2^{e_i-1}u_i)=1$ but $H(2^{e_i}u_i)=0$; hence the dyadic
flux is $1$ at arbitrarily large odd targets $u_{i+1}$.
Thus a bound on unweighted total flux per dyadic ray does not
control its size-weighted distribution.  The parity-sparse dangerous
language from the main manuscript is a separate restriction: the
present example disproves the proposed \emph{global odd-target}
implication, not a potentially stronger theorem localized to that
language.


\section{A positive refinement: dangerous odd targets have short initial valuations}
The failure of the \emph{global} odd-target condition does not rule out
using the sparse dangerous language from odd-core compression.  In
fact, that language automatically excludes very long initial
even runs.
\begin{lemma}[Valuation truncation in a dangerous cylinder]\label{lem:valuation}
Fix an odd $c\ge-1$ with $3\nmid c$, $1/2<\zeta_0<1$,
integers $2\le b\le t_*$, and let
$\mathcal C_{b,t_*}$ be the set of positive odd starting states $a$
under $T_{3,c}$ whose accelerated parity prefix of some length $t\in[b,t_*]$
contains at least $\zeta_0t$ odd states.  Then every
$a\in\mathcal C_{b,t_*}$ satisfies
\begin{equation}\label{eq:valuationcut}
 \boxed{\nu_2(3a+c)\le\left\lfloor(1-\zeta_0)t_*+1\right\rfloor.}
\end{equation}
In particular, for $\zeta_0=0.55$ only initial valuations
$e\le\lfloor0.45t_*+1\rfloor$ can contribute to the dangerous
odd-target potential sum.
\end{lemma}
\begin{proof}
Put $e=\nu_2(3a+c)\ge1$ and write $3a+c=2^eu$ with $u$ odd.
The accelerated orbit begins with the odd state $a$ and then has
$e-1$ consecutive even states before it reaches $u$ at time $e$.
Thus a prefix of length $t\le e$ contains exactly one odd state.
Since $t\ge b\ge2$ and $\zeta_0>1/2$, such a prefix cannot be
dangerous.  For a dangerous length $t>e$, the number of odd states
is at most $1+(t-e)$.  Therefore
$\zeta_0t\le1+t-e$, that is,
$e\le(1-\zeta_0)t+1\le(1-\zeta_0)t_*+1$.
\end{proof}
\begin{remark}[What this improves, and what it does not]
For the cutoff $t_*=t_X$ from our odd-core theorem, the relevant
\emph{dyadic flux exponents} are at most $0.45t_X+1$, rather than
$t_X$ itself.  This is a genuine deterministic localization,
not an additional assumption on an entrance law.  It does
\emph{not} shrink the modulus $2^{t_X}$ needed to recognize all
parity prefixes and does not by itself bound their weighted
predecessor mass.  On nonperiodic branching targets, that mass is
represented by differences $h(2^{e-1}u)-h(2^eu)$; periodic-target
corrections must be retained rather than silently omitted.
\end{remark}

\section{The genuine finite-source distinction}
The actual potential is the limit of logarithmically normalized
counting functions
\[
 h_N(a)=\frac1{\sum_{m\le N}1/m}
 \sum_{m\le N}\frac{\one_{\{\exists t\ge0:T^t(m)=a\}}}{m},
 \qquad h(a)=\lim_{N\to\infty}h_N(a).
\]
This origin as a limit is information \emph{not contained} in the
harmonic recurrence and cycle-level conservation.
\begin{proposition}[Finite-source vanishing along a selected ray]
For the sequence $(u_i)$ above and every fixed finite $N$,
there exists $i_N$ such that $h_N(u_i)=0$ for all $i\ge i_N$.
\end{proposition}
\begin{proof}
If the forward orbit of $m$ hits one of the $u_i$, it subsequently
reaches $1$ in finitely many steps.  Before reaching $1$ it has
only finitely many states, so it can hit only finitely many members
of the pairwise distinct sequence $(u_i)$.  Take the union of these
finite sets over $1\le m\le N$.
\end{proof}
\begin{remark}[Scope and next analytic input]
The theorem does \emph{not} disprove odd-core compression or show that
the actual $h$ fails odd-target uniform integrability.  It rules out
a strategy based only on boundedness, positivity, the harmonic
recurrence, dyadic telescoping, cycle-level conservation, or a bounded
first Ces\`aro moment.  The missing ingredient must exploit
quantitative \emph{distribution of logarithmic predecessor mass among
competing inverse branches}, or use the more localized dangerous
parity cylinders, without claiming that classical finite parity
bijections control later conditioned returns.  The proposition makes
clear why uniformity in the target $a$ matters when passing from
$h_N(a)$ to $h(a)$.
\end{remark}
\end{document}
